% 15. Источник: lecture_15_joint_distributions.tex: 7–28, 58–86. Страницы лекционного PDF: 37–38.
\question{Совместная функция распределения и независимость}

\begin{definition}[\topicmark{Случайный вектор}Случайный вектор]
    \textbf{Случайным вектором} называется вектор
    $\vec\xi=(\xi_1,\ldots,\xi_n)$, компоненты которого являются
    случайными величинами на одном вероятностном пространстве
    $(\Omega,\mathcal F,\Pr)$.
\end{definition}

\begin{definition}[\topicmark{Совместная функция распределения}Совместная функция распределения]
    \textbf{Функцией распределения случайного вектора} $\vec\xi$
    называется функция
    \[
        F_{\vec\xi}(x_1,\ldots,x_n)
        =\Pr\{\xi_1\leq x_1,\ldots,\xi_n\leq x_n\}.
    \]
    Иными словами,
    \[
        F_{\vec\xi}(\vec x)
        =\Pr\left(\bigcap_{i=1}^{n}\{\omega:\xi_i(\omega)\leq x_i\}\right).
    \]
    Она называется также \textbf{совместной функцией распределения}
    величин $\xi_1,\ldots,\xi_n$.
\end{definition}

\begin{definition}[\topicmark{Независимость через функции распределения}Определение]
    Случайные величины $\xi$ и $\eta$ называются \textbf{независимыми}, если
    \[
        F_{\xi,\eta}(x,y)=F_\xi(x)F_\eta(y)
        \qquad\forall x,y\in\R.
    \]
\end{definition}

\begin{theorem}[\topicmark{Независимость через события}Эквивалентное определение через события]
    Величины $\xi$ и $\eta$ независимы $\;\Leftrightarrow\;$
    для любых $A,B\in\mathcal B(\R)$ события $\{\xi\in A\}$ и $\{\eta\in B\}$
    независимы, то есть
    \[
        \Pr\{\xi\in A,\eta\in B\}
        =\Pr\{\xi\in A\}\Pr\{\eta\in B\}.
    \]
\end{theorem}

\begin{Proof}[От независимости событий к разложению функции распределения]
    Подставим $A=(-\infty;x]$, $B=(-\infty;y]$. Тогда
    \[
        \begin{aligned}
            F_{\xi,\eta}(x,y)
            &=\Pr\{\xi\leq x,\eta\leq y\}\\
            &=\Pr\{\xi\leq x\}\Pr\{\eta\leq y\}
             =F_\xi(x)F_\eta(y).
        \end{aligned}
    \]
\end{Proof}

% Обе требуемые формулировки и приведённый переход: Тервер(Лекция)67.pdf, с. 2–3.

% 16. Источник: lecture_15_joint_distributions.tex: 110–131. Страницы лекционного PDF: 39.
\question{Совместное дискретное распределение}

\begin{definition}[\topicmark{Совместное распределение}Определение]
    Пусть $\xi$ и $\eta$ принимают значения в конечных или счётных
    множествах $\{x_i\}$ и $\{y_j\}$ с вероятностью $1$.
    Их \textbf{совместное распределение} задаётся числами
    \[
        p_{ij}=\Pr\{\xi=x_i,\eta=y_j\}.
    \]
    Вероятности отдельных значений равны
    \[
        p_i=\Pr\{\xi=x_i\}=\sum_jp_{ij},\qquad
        q_j=\Pr\{\eta=y_j\}=\sum_ip_{ij}.
    \]
    При этом $p_{ij}\geq0$ и $\sum_{i,j}p_{ij}=1$.
\end{definition}

\begin{theorem}[\topicmark{Критерий независимости}Критерий независимости]
    Для дискретных случайных величин $\xi$ и $\eta$
    \[
        \xi\text{ и }\eta\text{ независимы}
        \quad\Leftrightarrow\quad
        p_{ij}=p_iq_j\quad\forall i,j.
    \]
\end{theorem}

% Тервер(Лекция)67.pdf, с. 2–4; Тервер(Лекция)69.pdf, с. 4 (факторизация вероятностей отдельных значений).
\begin{Proof}[$\Rightarrow$]
При независимости $\xi$ и $\eta$ события $\{\xi=x_i\}$ и $\{\eta=y_j\}$
независимы для каждой пары $i,j$. Поэтому
\[
 p_{ij}=\Pr\{\xi=x_i,\eta=y_j\}
       =\Pr\{\xi=x_i\}\Pr\{\eta=y_j\}=p_iq_j.
\]
\end{Proof}
\begin{Proof}[$\Leftarrow$\editorialmark]
Пусть $p_{ij}=p_iq_j$ для всех $i,j$. Для любых борелевских множеств
$A,B\subseteq\R$ событие $\{\xi\in A,\eta\in B\}$ объединяет попарно
несовместные исходы таблицы с $x_i\in A$ и $y_j\in B$. Поэтому
\[
 \begin{aligned}
 \Pr\{\xi\in A,\eta\in B\}
 &=\sum_{i:\,x_i\in A}\sum_{j:\,y_j\in B} p_iq_j\\
 &=\left(\sum_{i:\,x_i\in A}p_i\right)
   \left(\sum_{j:\,y_j\in B}q_j\right)
 =\Pr\{\xi\in A\}\Pr\{\eta\in B\}.
 \end{aligned}
\]
Получили определение независимости через события.
\end{Proof}

\begin{remark}[Матричная запись для конечных распределений]
Для конечной таблицы $P=(p_{ij})$ критерий записывается также в виде
$\xi$ и $\eta$ независимы $\Leftrightarrow \rk P=1$.
При независимости $i$-я строка таблицы равна $p_i(q_1,\ldots,q_s)$.
\end{remark}

\begin{theorem}[\topicmark{Свёртка}Дискретная свёртка\editorialmark]
Пусть $\xi$ и $\eta$ независимы и дискретны. Чтобы найти вероятность
значения $z$ у суммы, складываем вероятности всех пар, дающих $z$:
\[
 \Pr\{\xi+\eta=z\}=\sum_{(i,j):\,x_i+y_j=z}p_iq_j.
\]
Для целочисленных величин стоит зафиксировать первое слагаемое $k$.
Тогда второе должно равняться $m-k$, и при $m\in\Z$
\[
 \Pr\{\xi+\eta=m\}
 =\sum_{k\in\Z}\Pr\{\xi=k\}\Pr\{\eta=m-k\}.
\]
\end{theorem}

% 17. Источник: lecture_15_joint_distributions.tex: 156–205. Страницы лекционного PDF: 40.
\question{Совместное абсолютно непрерывное распределение}

\begin{definition}[\topicmark{Совместная плотность}Определение]
    Распределение пары $(\xi,\eta)$ называется \textbf{абсолютно непрерывным},
    если существует неотрицательная измеримая функция $p_{\xi,\eta}$,
    для которой
    \[
        F_{\xi,\eta}(x,y)
        =\int_{-\infty}^{y}\int_{-\infty}^{x}
            p_{\xi,\eta}(u,v)\,\diff u\,\diff v.
    \]
    Функция $p_{\xi,\eta}$ называется \textbf{совместной плотностью}.
    Она удовлетворяет условию
    \[
        \iint_{\R^2}p_{\xi,\eta}(u,v)\,\diff u\,\diff v=1.
    \]
\end{definition}

\begin{theorem}[\topicmark{Критерий независимости}Критерий независимости через плотности]
    Если пара $(\xi,\eta)$ имеет совместную плотность, то
    \[
        \xi\text{ и }\eta\text{ независимы}
        \quad\Longleftrightarrow\quad
        p_{\xi,\eta}(x,y)=p_\xi(x)p_\eta(y)
        \ \text{почти всюду}.
    \]
    Здесь почти всюду понимается относительно меры Лебега на $\R^2$.
\end{theorem}

\begin{Proof}[\topicmark{\texorpdfstring{$\Rightarrow$}{⇒}}$\Rightarrow$]
    При независимости $F_{\xi,\eta}(x,y)=F_\xi(x)F_\eta(y)$.
    Дифференцируя по обеим переменным, получаем почти всюду
    \[
        p_{\xi,\eta}(x,y)
        =\frac{\partial^2F_{\xi,\eta}}{\partial x\,\partial y}(x,y)
        =F_\xi'(x)F_\eta'(y)=p_\xi(x)p_\eta(y).
    \]
\end{Proof}

\begin{Proof}[\topicmark{\texorpdfstring{$\Leftarrow$}{⇐}}$\Leftarrow$]
    Если $p_{\xi,\eta}(u,v)=p_\xi(u)p_\eta(v)$ почти всюду, то
    \[
        \begin{aligned}
            F_{\xi,\eta}(x,y)
            &=\int_{-\infty}^{y}\int_{-\infty}^{x}
                p_\xi(u)p_\eta(v)\,\diff u\,\diff v\\
            &=\left(\int_{-\infty}^{x}p_\xi(u)\,\diff u\right)
              \left(\int_{-\infty}^{y}p_\eta(v)\,\diff v\right)\\
            &=F_\xi(x)F_\eta(y).
        \end{aligned}
    \]
\end{Proof}

\begin{theorem}[\topicmark{Свёртка}Свёртка плотностей\editorialmark]
Пусть $\xi$ и $\eta$ независимы и имеют плотности $p_\xi$ и $p_\eta$.
Тогда $S=\xi+\eta$ имеет плотность, которая почти всюду равна
\[
    p_S(z)=\int_{\R}p_\xi(u)\,p_\eta(z-u)\,\diff u.
\]
При фиксированном $z$ вклад дают те $u$, для которых обе плотности
в подынтегральном произведении отличны от нуля.
\end{theorem}
